\documentclass[11pt]{article}
\usepackage[margin=1in]{geometry}
\usepackage{amsmath, amssymb}
\usepackage{booktabs}
\usepackage{setspace}
\usepackage{enumitem}

\setstretch{1.2}
\setlength{\parindent}{0pt}
\setlist[itemize]{itemsep=0.4em}

\begin{document}

\begin{center}
{\Large \textbf{Tutorial: Inequality and Poverty in Canada}}\\
\vspace{1em}
{\large ECON 204 -- Contemporary Economic Issues}\\
\vspace{1em}
{\large Winter 2026}
\end{center}
\hrule
\vspace{1em}


%----------------------------------------------------
\section*{Question 1. Income Distribution and Inequality}

Consider the following \textbf{annual before-tax household incomes} (in Canadian dollars) for \textbf{10 households}, ordered from lowest to highest income:

\begin{center}
\begin{tabular}{cc}
\toprule
\textbf{Household} & \textbf{Income (\$)} \\
\midrule
1 & 18{,}000 \\
2 & 22{,}000 \\
3 & 26{,}000 \\
4 & 30{,}000 \\
5 & 35{,}000 \\
6 & 42{,}000 \\
7 & 50{,}000 \\
8 & 65{,}000 \\
9 & 95{,}000 \\
10 & 180{,}000 \\
\bottomrule
\end{tabular}
\end{center}

\subsection*{(a) Mean and Median Income}

\begin{itemize}
\item Calculate the \textbf{mean (average)} household income.
\item Calculate the \textbf{median} household income.
\end{itemize}

Explain which measure better reflects the ``typical'' household income in Canada and why.

\subsection*{(b) Income Shares}

\begin{itemize}
\item Calculate total income across all households.
\item Calculate the income share of:
\begin{itemize}
\item the \textbf{bottom 40\%} (households 1--4),
\item the \textbf{top 10\%} (household 10).
\end{itemize}
\end{itemize}

Use the formula:
\[
\text{Income Share} = \frac{\text{Group Income}}{\text{Total Income}} \times 100
\]

Interpret what these shares imply about income concentration.

%----------------------------------------------------
\section*{Question 2. Simple Inequality Ratios}

\subsection*{(a) Top-to-Bottom Ratio}

Calculate the ratio:
\[
\text{Top-to-Bottom Ratio} = \frac{\text{Income of richest household}}{\text{Income of poorest household}}
\]

Explain what this ratio indicates about inequality.

\subsection*{(b) Policy Change}

Suppose a progressive tax-and-transfer policy reduces the richest household's income to \$150{,}000 and raises the poorest household's income to \$25{,}000.

\begin{itemize}
\item Recalculate the top-to-bottom ratio.
\item State whether inequality has increased or decreased.
\end{itemize}

%----------------------------------------------------
\section*{Question 3. Poverty Measurement in Canada}

Assume the following annual poverty thresholds for a \textbf{single adult living in a large Canadian city}:

\begin{itemize}
\item Low-Income Cut-Off (LICO): \$26{,}000
\item Market Basket Measure (MBM): \$28{,}000
\end{itemize}

\subsection*{(a) Poverty Status}

Using the income data from Question 1:
\begin{itemize}
\item How many households fall below the LICO?
\item How many households fall below the MBM?
\end{itemize}

\subsection*{(b) Poverty Rates}

Calculate the poverty rate under each measure:
\[
\text{Poverty Rate} = \frac{\text{Number of households below the threshold}}{\text{Total households}} \times 100
\]

Explain why the poverty rate differs across measures.

%----------------------------------------------------
\section*{Question 4. Depth of Poverty}

For households below the MBM threshold:

\subsection*{(a) Poverty Gap}

For each poor household, calculate:
\[
\text{Poverty Gap}_i = \text{MBM Threshold} - \text{Household Income}_i
\]

\subsection*{(b) Average Poverty Gap}

Compute the average poverty gap among households below the MBM.

Explain why the poverty gap is an important policy-relevant measure.

%----------------------------------------------------
\section*{Question 5. Inflation and Poverty}

Suppose inflation increases the MBM threshold by \textbf{8\%}, while nominal household incomes remain unchanged.

\subsection*{(a) New Poverty Line}

Calculate the new MBM threshold.

\subsection*{(b) Poverty Impact}

\begin{itemize}
\item How many households now fall below the new MBM threshold?
\item Explain how inflation can increase poverty even if nominal incomes do not change.
\end{itemize}



\end{document}
